\documentclass[11pt,a4paper]{article}

\usepackage[utf8]{inputenc}

\usepackage[T1]{fontenc}

\usepackage[polish]{babel}

\usepackage{lmodern}

\usepackage{geometry}

\geometry{margin=27mm}

\usepackage[hidelinks]{hyperref}

\title{Neuronaukowe wyjaśnianie krzywdy i granice wnioskowania odwrotnego}

\author{}

\date{Szkic naukowy, wersja 0.1.1}

\begin{document}

\maketitle

\begin{abstract}

Tematem pracy jest Neuronaukowe wyjaśnianie krzywdy i granice wnioskowania odwrotnego. Tekst systematyzuje uzasadnienie, hipotezy i sposób weryfikacji zagadnienia. Dobór pierwszorzędowego wskaźnika pozostaje elementem przyszłego protokołu. Nie przedstawia się niezweryfikowanych wyników ani fikcyjnych cytowań.

\end{abstract}

\section{Wprowadzenie}

Postawy społeczne, przywiązanie i normy moralne wymagają rozróżnienia deklaracji, zachowania i sytuacji społecznej. Zagadnienie Neuronaukowe wyjaśnianie krzywdy i granice wnioskowania odwrotnego wymaga odróżnienia badanego konstruktu od jego operacjonalizacji i kontekstowych źródeł zmienności. Trzeba porównać możliwe wyjaśnienia i ustalić granice interpretacji dowodów.

\section{Cel i pytania badawcze}

Celem pracy jest rozstrzygnięcie zagadnienia: Neuronaukowe wyjaśnianie krzywdy i granice wnioskowania odwrotnego. Dobór pierwszorzędowego wskaźnika pozostaje elementem przyszłego protokołu. Hipotezy i ich kierunki muszą zostać określone przed analizą danych lub krytyczną syntezą źródeł.

\section{Materiał i metody}

Zakodować opisy neuronaukowych korelatów krzywdy i odróżnić aktywację mózgową od swoistości procesu. Pierwszorzędowy wskaźnik: Częstość nieuprawnionego wnioskowania odwrotnego w analizowanych publikacjach. Nie przedstawia się przeprowadzonych analiz ani uzyskanych wyników. Przed realizacją konieczne jest określenie źródeł materiału, kryteriów jakości, zasad reprodukowalności i licencji.

\section{Plan analizy}

Proponowany plan analityczny: podwójne kodowanie, zgodność ekspertów i analiza typów błędów. Zamrożenie estymandu, kryteriów jakości i analiz wrażliwości nastąpi przed dostępem do danych; na tym etapie nie ma oszacowań ani wyników.

\section{Wyniki}

W niniejszym pliku nie umieszczono wyników empirycznych, statystyk ani nieprzeprowadzonych analiz. Sekcja wymaga uzupełnienia po uzyskaniu i weryfikacji materiału dowodowego.

\section{Dyskusja i ograniczenia}

Ograniczenia obejmują możliwe błędy pomiaru, dobór próby, alternatywne mechanizmy i ograniczoną możliwość uogólniania. Korelacja nie dowodzi przyczynowości; poprawność techniczna nie jest dowodem znaczenia klinicznego.

\section{Wnioski}

Zdefiniowano ramy opracowania, ale nie sformułowano wniosków wykraczających poza dostępny stan dowodów.

\section*{Bibliografia}

Weryfikacja literatury i przypisanie cytowań do stwierdzeń są kolejnym etapem prac.

\end{document}

